\documentclass[11pt]{article}
\usepackage[margin=1in]{geometry}
\usepackage{amsmath}
\usepackage{amssymb}
\usepackage{enumitem}
\usepackage{xcolor}
\usepackage{fancyhdr}
\usepackage{listings}
\usepackage[colorlinks=true, linkcolor=blue, urlcolor=blue]{hyperref}

% Define OSU orange and AI blue
\definecolor{osaorange}{RGB}{220,115,38}

% Header/Footer
\pagestyle{fancy}
\fancyhf{}
\lhead{H524 Introduction to Biostatistics}
\rhead{Week 4 Assignment}
\cfoot{\thepage}

% Answer box styling
\newcommand{\answer}[1]{\par\vspace{0.3em}\colorbox{green!10}{\parbox{\dimexpr\textwidth-2\fboxsep}{\textbf{ANSWER:} #1}}\vspace{0.3em}\par}

\title{\textbf{\Large H524 Week 4 Assignment: Confidence Intervals}\\
\large Student Version}
\author{Introduction to Biostatistics}
\date{Fall 2025}

\begin{document}

\maketitle

\begin{center}
\fbox{\parbox{0.9\textwidth}{
\textbf{Assignment Information:}
\begin{itemize}[leftmargin=*]
\item \textbf{Total Points:} 110 points
\item \textbf{Time Limit:} 180 minutes (3 hours)
\item \textbf{Attempts:} 1
\item \textbf{Question Types:} 16 Multiple Choice (80 pts), 2 Essay/Fill-in (30 pts)
\end{itemize}
}}
\end{center}

\vspace{1em}

\section*{Part 1: CI Fundamentals}

\subsection*{Question 1: CI Interpretation (5 points)}

A study reports a 95\% confidence interval for mean blood pressure reduction as (8.2, 14.6) mmHg. What is the CORRECT interpretation?

\begin{enumerate}[label=\alph*)]
\item We are 95\% confident that the true population mean blood pressure reduction is between 8.2 and 14.6 mmHg
\item There is a 95\% probability that the true mean is between 8.2 and 14.6 mmHg
\item 95\% of patients will have blood pressure reductions between 8.2 and 14.6 mmHg
\item The sample mean has a 95\% chance of being between 8.2 and 14.6 mmHg
\end{enumerate}

\subsection*{Question 2: Critical Value Selection (5 points)}

You want to construct a 95\% confidence interval for a population mean. The sample size is $n = 15$, and the population standard deviation is unknown. Which critical value should you use?

\begin{enumerate}[label=\alph*)]
\item $z = 1.96$
\item $t(14) = 2.145$
\item $t(15) = 2.131$
\item $z = 2.576$
\end{enumerate}

\subsection*{Question 3: Factors Affecting CI Width (5 points)}

If you want a NARROWER confidence interval for a population mean, which of the following actions would achieve this?

\begin{enumerate}[label=\alph*)]
\item Increase the confidence level from 95\% to 99\%
\item Decrease the sample size
\item Increase the sample size
\item Use a larger standard deviation
\end{enumerate}

\subsection*{Question 4: Standard Error (5 points)}

A study has a sample standard deviation $s = 12$ and sample size $n = 36$. What is the standard error of the mean?

\begin{enumerate}[label=\alph*)]
\item $SE = 2$
\item $SE = 3$
\item $SE = 6$
\item $SE = 12$
\end{enumerate}

\section*{Part 2: CI for Means}

\subsection*{Question 5: Manual CI Calculation (15 points)}

A clinical trial measures pain reduction scores for 25 patients receiving a new treatment. The results are:
\begin{itemize}
\item Sample mean: $\bar{x} = 7.8$ points
\item Sample standard deviation: $s = 3.2$ points
\item Sample size: $n = 25$
\end{itemize}

Calculate a 95\% confidence interval for the true mean pain reduction. Show all work including:
\begin{enumerate}[label=\alph*)]
\item Which distribution to use ($z$ or $t$) and why (3 points)
\item The appropriate critical value (3 points)
\item The standard error calculation (3 points)
\item The margin of error calculation (3 points)
\item The final confidence interval with interpretation (3 points)
\end{enumerate}

\textit{Hint: For a 95\% CI with df=24, the critical $t$-value is approximately 2.064}

\subsection*{Question 6: Confidence Level Comparison (5 points)}

Using the same data from the Manual CI question above, if we constructed a 99\% confidence interval instead of 95\%, what would happen to the interval width?

\begin{enumerate}[label=\alph*)]
\item The interval would be narrower because we are more confident
\item The interval would be wider because higher confidence requires a larger margin of error
\item The interval width would stay exactly the same
\item The interval would be narrower because the critical value decreases
\end{enumerate}

\section*{Part 3: CI for Proportions}

\subsection*{Question 7: Proportion CI Conditions (5 points)}

When constructing a confidence interval for a population proportion using the normal approximation (Wald method), which condition must be satisfied?

\begin{enumerate}[label=\alph*)]
\item $n \geq 30$
\item Both $np \geq 5$ and $n(1-p) \geq 5$
\item The population must be normally distributed
\item The sample proportion must be greater than 0.5
\end{enumerate}

\subsection*{Question 8: Proportion CI Calculation (5 points)}

In a survey of 200 patients, 150 reported satisfaction with their treatment. What is the sample proportion ($\hat{p}$)?

\begin{enumerate}[label=\alph*)]
\item $\hat{p} = 0.50$
\item $\hat{p} = 0.67$
\item $\hat{p} = 0.75$
\item $\hat{p} = 0.80$
\end{enumerate}

\subsection*{Question 9: Standard Error for Proportion (5 points)}

Using the data from the proportion calculation question ($\hat{p} = 0.75$, $n = 200$), what is the standard error for the proportion? Use $SE = \sqrt{\frac{\hat{p}(1-\hat{p})}{n}}$

\begin{enumerate}[label=\alph*)]
\item $SE \approx 0.031$
\item $SE \approx 0.043$
\item $SE \approx 0.075$
\item $SE \approx 0.125$
\end{enumerate}

\subsection*{Question 10: Vaccine Efficacy CI (5 points)}

A vaccine trial reports a 95\% CI for vaccine efficacy as (72\%, 89\%). What can we conclude?

\begin{enumerate}[label=\alph*)]
\item We are 95\% confident the true efficacy is between 72\% and 89\%, suggesting the vaccine is effective since the entire interval is above 0\%
\item 72\% of people are protected at minimum, 89\% at maximum
\item The vaccine has a 95\% chance of working
\item We cannot determine effectiveness because the interval is too wide
\end{enumerate}

\section*{Part 4: Sample Size Determination}

\subsection*{Question 11: Sample Size Concepts (5 points)}

You want to estimate the mean cholesterol level in a population with a margin of error of 5 mg/dL at 95\% confidence. If you quadruple the sample size, what happens to the margin of error?

\begin{enumerate}[label=\alph*)]
\item The margin of error is cut in half
\item The margin of error is cut to one-quarter
\item The margin of error doubles
\item The margin of error stays the same
\end{enumerate}

\subsection*{Question 12: Sample Size Calculation (5 points)}

A researcher wants to estimate a population proportion with a 95\% CI and margin of error of 0.05. Using the conservative estimate ($\hat{p} = 0.5$), what sample size is needed?

Formula: $n = \frac{z^2 \cdot \hat{p}(1-\hat{p})}{E^2}$, where $z = 1.96$ for 95\% CI

\begin{enumerate}[label=\alph*)]
\item $n = 196$
\item $n = 271$
\item $n = 385$
\item $n = 577$
\end{enumerate}

\section*{Part 5: Assumptions and Checking}

\subsection*{Question 13: Normality Assessment (5 points)}

You have a sample of $n = 12$ observations and want to construct a $t$-based confidence interval for the mean. Which of the following is the BEST way to assess whether the normality assumption is reasonable?

\begin{enumerate}[label=\alph*)]
\item Check that $n \geq 30$, which automatically ensures normality
\item Create a histogram and Q-Q plot to visually assess normality, and consider the Shapiro-Wilk test
\item Calculate the mean and median; if they're close, the data is normal
\item Normality is not required for $t$-based intervals
\end{enumerate}

\subsection*{Question 14: Large vs Small Sample (5 points)}

When can you safely use a $z$-based confidence interval instead of a $t$-based confidence interval for a population mean?

\begin{enumerate}[label=\alph*)]
\item When the population standard deviation ($\sigma$) is known, regardless of sample size
\item Only when $n \geq 30$
\item When the sample size is less than 30
\item Never; always use $t$-distribution for real-world data
\end{enumerate}

\subsection*{Question 15: CI for Hypothesis Testing (5 points)}

A 95\% confidence interval for the difference in mean blood pressure between two treatments is $(-2.3, 8.7)$ mmHg. What can you conclude at the $\alpha = 0.05$ significance level?

\begin{enumerate}[label=\alph*)]
\item There is a significant difference between treatments because the interval is positive
\item There is NOT a significant difference between treatments because the interval contains 0
\item Treatment 1 is better than Treatment 2
\item We need more data to make any conclusion
\end{enumerate}

\section*{Part 6: R Implementation}

\subsection*{Question 16: CI with R (15 points)}

A clinical trial measures vitamin D levels (ng/mL) for 30 patients after supplementation:

\begin{lstlisting}[language=R, basicstyle=\small\ttfamily]
vitamin_d <- c(32, 38, 29, 35, 31, 40, 27, 34, 36, 30,
               28, 37, 33, 35, 31, 34, 36, 29, 33, 37,
               30, 35, 34, 31, 39, 33, 36, 29, 34, 35)
\end{lstlisting}

Using R, complete the following tasks and paste your code and results:
\begin{enumerate}[label=\alph*)]
\item Calculate the sample mean, standard deviation, and sample size (3 points)
\item Create a histogram to assess normality (2 points)
\item Use \texttt{t.test()} to calculate a 95\% confidence interval for the mean vitamin D level (5 points)
\item Interpret the confidence interval in context (3 points)
\item What would you conclude about whether the mean vitamin D level is above 30 ng/mL? (2 points)
\end{enumerate}

\textbf{a) Descriptive statistics (3 points):}
\begin{lstlisting}[language=R, basicstyle=\small\ttfamily]
vitamin_d <- c(32, 38, 29, 35, 31, 40, 27, 34, 36, 30,
               28, 37, 33, 35, 31, 34, 36, 29, 33, 37,
               30, 35, 34, 31, 39, 33, 36, 29, 34, 35)

# Sample statistics
mean(vitamin_d)
# [1] 33.37

sd(vitamin_d)
# [1] 3.35

length(vitamin_d)
# [1] 30
\end{lstlisting}

\textbf{b) Histogram for normality (2 points):}
\begin{lstlisting}[language=R, basicstyle=\small\ttfamily]
hist(vitamin_d,
     main = "Histogram of Vitamin D Levels",
     xlab = "Vitamin D (ng/mL)",
     col = "lightblue",
     breaks = 10)
\end{lstlisting}
The histogram should show approximately normal distribution.

\textbf{c) 95\% CI using t.test() (5 points):}
\begin{lstlisting}[language=R, basicstyle=\small\ttfamily]
t.test(vitamin_d, conf.level = 0.95)

# Output:
# One Sample t-test
#
# data:  vitamin_d
# t = 54.595, df = 29, p-value < 2.2e-16
# alternative hypothesis: true mean is not equal to 0
# 95 percent confidence interval:
#  32.11669 34.61664
# sample estimates:
# mean of x
#    33.36667
\end{lstlisting}

\textbf{d) Interpretation (3 points):}

We are 95\% confident that the true mean vitamin D level for all patients receiving this supplementation is between 32.12 and 34.62 ng/mL.

\textbf{e) Conclusion about 30 ng/mL (2 points):}

Since the entire 95\% confidence interval (32.12, 34.62) is above 30 ng/mL, we can conclude with 95\% confidence that the true mean vitamin D level is greater than 30 ng/mL. This suggests the supplementation successfully raises vitamin D levels above the 30 ng/mL threshold.

\subsection*{Question 17: Proportion CI in R (5 points)}

Which R function would you use to calculate a confidence interval for a population proportion?

\begin{enumerate}[label=\alph*)]
\item \texttt{prop.test(x, n)}
\item \texttt{t.test(x)}
\item \texttt{binom.test(x, n)}
\item Both a and c are appropriate
\end{enumerate}

\end{document}